\section*{9. 矩阵连乘问题}

给定$n$个矩阵$M_1, M_2, \ldots, M_n$，其中$M_i$维数为$r_{i-1} \times r_i$。矩阵乘法满足结合律，不同的括号化方式导致不同的乘法次数。求使乘法运算次数最少的最优括号化方案。



\begin{itemize}
    \item \textbf{状态定义}：$m[i][j]$表示矩阵$M_i$到$M_j$连乘的最小乘法次数
    \item \textbf{转移方程}：$$m[i][j] = \min_{i \le k < j} \{m[i][k] + m[k+1][j] + r_{i-1} \times r_k \times r_j\}$$
    \item \textbf{边界条件}：$m[i][i] = 0$
    \item \textbf{计算顺序}：按区间长度$x = j-i+1$从$1$到$n$递增
    \item \textbf{时间复杂度}：$O(n^3)$
\end{itemize}
\begin{lstlisting}[language=C++, style=cppstyle]
int matrixChainOrder(vector<int>& r) {
    int n = r.size() - 1;  // n个矩阵
    vector<vector<int>> m(n + 1, vector<int>(n + 1, 0));

    // 按区间长度递增
    for (int x = 2; x <= n; x++) {  // x = j-i+1
        for (int i = 1; i <= n - x + 1; i++) {
            int j = i + x - 1;
            m[i][j] = INT_MAX;

            for (int k = i; k < j; k++) {
                int cost = m[i][k] + m[k+1][j] + r[i-1] * r[k] * r[j];
                if (cost < m[i][j]) {
                    m[i][j] = cost;
                }
            }
        }
    }

    return m[1][n];
}

// 带最优解构造版本
pair<int, string> matrixChainWithSolution(vector<int>& r) {
    int n = r.size() - 1;
    vector<vector<int>> m(n + 1, vector<int>(n + 1, 0));
    vector<vector<int>> s(n + 1, vector<int>(n + 1, 0));  // 记录切分点

    for (int x = 2; x <= n; x++) {
        for (int i = 1; i <= n - x + 1; i++) {
            int j = i + x - 1;
            m[i][j] = INT_MAX;

            for (int k = i; k < j; k++) {
                int cost = m[i][k] + m[k+1][j] + r[i-1] * r[k] * r[j];
                if (cost < m[i][j]) {
                    m[i][j] = cost;
                    s[i][j] = k;
                }
            }
        }
    }

    return {m[1][n], "可通过s数组回溯构造括号化方案"};
}
\end{lstlisting}


\begin{enumerate}
    \item 建二维表，m[i][j]表示从第i个矩阵乘到第j个矩阵的最少乘法次数
    \item 按区间长度从小到大算
    \item 枚举分割点k：代价等于"左半边"加"右半边"加"最后一次乘法"
    \item 取所有k中代价最小的那个
\end{enumerate}

